
\section{Recommendations}
\label{sec:recommendations_bench}

Based on the benchmark evaluation, we provide the following guidelines:

\subsection{Full-Model vs.\ Residual Learning}

\textbf{Full-model learning} (NN, SINDy, DMD, GP, Hankel) should be used when
no analytical model is available, or when interpretability, linear MPC
compatibility, or uncertainty quantification are required. Among these, SINDy
provides the best disturbed-condition accuracy while remaining physically
interpretable and fast enough for real-time MPC.

\textbf{Residual learning} (PINN) is preferred when a partial analytical model
exists and the expected residual is small. It achieves near-zero prediction
error on clean data by design, and requires fewer samples.

\subsection{When to Use Both}

The \textit{online adaptation gap} of \textbf{0.0004~m}
quantifies the degradation of the best residual model under disturbances.
This gap cannot be closed by better offline training; it demands \textit{online
adaptation}. We recommend the following pipeline:
\begin{enumerate}
  \item Train a PINN residual model offline on clean data.
  \item Deploy with a LoRA adapter initialised at zero (no residual correction).
  \item Update the LoRA adapter online using Experience Replay or EWC to
        compensate for observed disturbances without forgetting clean-water behaviour.
\end{enumerate}
This combination achieves the nominal accuracy of PINN and the disturbance
rejection of LoRA continual learning, as demonstrated in the companion paper.
